% =================================================================
% CONFIGURAÇÕES DO TEMPLATE - PAPER VARIANT
% Variant streamlined for 6-10 page conference/journal submissions.
% Como usar:
%  - Preencha os placeholders abaixo (título, autor, etc.)
%  - Adicione/remova pacotes conforme a necessidade do seu trabalho
%  - Mantenha overrides aqui, sem editar config/sbc-template.sty
% =================================================================

% Pacotes essenciais
\usepackage{config/sbc-template}
\usepackage[T1]{fontenc}   % melhor hifenização/acentuação com pdflatex
\usepackage[utf8]{inputenc}
\usepackage{microtype}     % microtipografia (kerning/hifenização)
\usepackage[english,brazil]{babel} % idiomas: en + pt-BR
\usepackage{csquotes}      % citações tipográficas
\usepackage{graphicx,url}
\graphicspath{{assets/}}   % pasta padrão para imagens

% Pacotes adicionais (essenciais para papers)
\usepackage{listings}           % listagens de código
\usepackage{xcolor}             % cores (necessário para listings)
\usepackage{amsmath,amssymb}    % matemática
\usepackage{booktabs}           % tabelas profissionais (\toprule, etc.)
\usepackage{hyperref}           % hyperlinks

% Estilo de código (exemplo; ajuste ou remova)
\definecolor{codegray}{RGB}{248,248,248}
\lstdefinestyle{code}{
    backgroundcolor=\color{codegray},
    basicstyle=\footnotesize\ttfamily,
    numbers=left,
    numbersep=5pt,
    frame=single,
}

% Hyperlinks (preenche metadados do PDF)
\hypersetup{
    colorlinks=true,
    linkcolor=black,
    urlcolor=blue,
    citecolor=blue,
    pdftitle={Test Project},
    pdfauthor={Test Author},
}

% Comandos de placeholders (consumidos em \title, \author e \address)
\newcommand{\DocumentTitle}{Test Project}
\newcommand{\AuthorName}{Test Author}
\newcommand{\AuthorInstitution}{Test Uni}
\newcommand{\AuthorLocation}{Test City}
\newcommand{\AuthorEmail}{test@example.com}

% Metadados padrão (para projetos que usam este arquivo como entrada principal)
\title{\DocumentTitle}
\author{\AuthorName\inst{1}}
\address{\AuthorInstitution\\
  \AuthorLocation\\
  \email{\AuthorEmail}
}

% Configurações gerais
\sloppy
\setlength{\parskip}{6pt}
